\enlargethispage{2\baselineskip}%
% Abstract (NBR 6028: one paragraph, no formulas).
For a decade, model selection for tabular data relied on the maxim that
gradient boosting always wins. Tabular foundation models, pre-trained
networks that predict without task-specific training, have cast doubt on it.
However, most available evidence comes from the proponents themselves, under
differing protocols. This dissertation produces evidence under a single
protocol and turns it into a selection procedure, along three lines. First,
a multi-criteria benchmark evaluates 14 models from three families
(gradient boosting, deep learning, and foundation models) on 18 OpenML
datasets, with a hold-out untouched by selection, frequentist and
Bayesian inference, and measured training, tuning, and latency costs.
Second, a decision framework is assessed: its learned component under
leave-one-dataset-out validation, its fixed policies descriptively. Third, the distillation of the TabPFN teacher's quantile curves into
gradient-boosting students is studied on 20 regression datasets, under registered
hypotheses. Among the 11 core models (prior to the 2026 extension), the top
of the accuracy ranking, measured in binary classification, is a tie. Other
tasks are read descriptively. Under a ROPE of one AUC percentage point,
3 of the other 10 models are practically equivalent to the leader, TabPFN, which is itself a foundation model,
by the strict Bayesian criterion, and 8 by the lenient one. In the 2026 generation, the untuned TabFM has the best average rank in all three tasks (in multiclass, over the two
datasets where it operates) and is the best model on 13 of 18 datasets
(operating on 17). The lead is consistent across datasets, but the
underpowered tests do not confirm it as a statistical separation, and this lead costs
about five orders of magnitude more latency on the reference dataset
(\texttt{adult}). Kolmogorov--Arnold networks finish in the bottom third on
all three tasks. For the decision, a router learned from meta-features is
less accurate than the trivial rule (0.11 vs.\ 0.44). The fixed
``foundation model first'' policy (an intra-family oracle) has zero median
regret on the primary metric, but stops being the recommendation
once 10 to 20 percent of the decision weight shifts from accuracy to
latency, by median regret, under rank normalization. In distillation, the
student trained on the teacher's means beats the control on only 4 of 20
datasets. The student trained on the quantile curves retains, at the
median, 12 to 19 percent of the distributional advantage, depending on how the arms are tuned, and retention is positive on 16 of 20 (sign test, p
approximately 0.019 after the sequential correction). It beats
default-configured classical distributional baselines on 17 of 20 (12 of 20
against the per-dataset best alternative). In the regime studied (up to 100
thousand samples, pinned versions, a single GPU), model choice should start
not from a universal winner but from the application's constraints:
latency, architectural limits, and auditable explanations. In this regime,
distillation recovers a modest median share of the foundation model's
predictive uncertainty, under preconditions measurable before training the student. The advantage gate was verified only in an oracle version, and
the rule was assessed only on the datasets it was derived from.

\vspace{0.4cm}
\noindent\textbf{Keywords:} tabular data; foundation models;
benchmarking; knowledge distillation; distributional regression.
